
RailFlux Route Scanning Procedure - Technical Analysis
Procedure Initialization and Entry Point
Scan Invocation Framework: Route scanning commences through scanDestinationSignals method accepting sourceSignalId parameter, direction specification (UP/DOWN/AUTO), and includeBlocked boolean flag determining whether blocked destinations appear in results. Performance tracking begins immediately with QElapsedTimer establishing scan duration metrics for system monitoring and optimization analysis.
Direction Determination Logic: AUTO direction parameter triggers determineSignalDirection method extracting direction attribute from database signal record, defaulting to UP when database connectivity fails or signal lacks direction specification. Manual direction specification bypasses determination logic while requiring validation against UP/DOWN enumeration constraints before proceeding with candidate evaluation.
Source Signal Validation: Database connectivity verification precedes signal existence validation through getSignalById query returning empty QVariantMap for nonexistent signals. Source signal validation failure terminates scanning immediately with descriptive error message indicating specific failure condition preventing scan completion.
Signal Type Compatibility Matrix Implementation
Hierarchical Signal Progression Rules: Signal type compatibility matrix implements railway operational hierarchy through QMap data structure defining permitted source-to-destination signal progressions. HOME signals connect exclusively to STARTER destinations, STARTER signals progress to ADVANCED_STARTER targets, ADVANCED_STARTER signals maintain extensible destination arrays for future enhancement, OUTER signals optionally connect to HOME destinations establishing optional approach-control relationships.
Database Query Construction: Eligible destination identification constructs PostgreSQL queries against v_signals_complete view incorporating direction matching, activation status verification, route_signal capability requirements, manual_control exclusion, circuit connectivity validation through preceded_by_circuit_id and succeeded_by_circuit_id constraints, and source signal exclusion preventing self-referential route requests.
PostgreSQL Array Format Conversion: QStringList compatibility array transforms into PostgreSQL array syntax through join operations with curly brace encapsulation, enabling ANY operator matching within database query execution. Array parameter binding utilizes prepared statement placeholders ensuring SQL injection protection while maintaining query performance through statement caching.
Individual Candidate Evaluation Pipeline
Destination Candidate Structure: DestinationCandidate objects encapsulate comprehensive evaluation results including destSignalId reference, direction specification, displayName composed from signal name and type information, reachability enumeration (REACHABLE_CLEAR/REACHABLE_REQUIRES_PM/BLOCKED), blockedReason textual explanations, pathSummary containing hop count and weight metrics, requiredPMActions arrays specifying necessary point machine movements, and conflicts arrays identifying specific blocking conditions.
Circuit Resolution Process: Signal-to-circuit mapping extracts succeededByCircuitId from source signal and precededByCircuitId from destination signal establishing pathfinding endpoint parameters. Circuit resolution failure triggers INCOMPLETE_TOPOLOGY blocking with candidate termination, while successful resolution enables GraphService pathfinding invocation with current point machine states and configurable timeout parameters.
GraphService Integration: A-star pathfinding algorithm execution through GraphService.findRoute method processes start circuit, goal circuit, direction constraints, point machine state mappings, and timeout limitations returning comprehensive pathfinding results including success status, discovered path as circuit identifier array, cost calculations, processing time metrics, and explored node counts for performance analysis.
Path Analysis and Clearance Validation
Path Metrics Computation: Successful pathfinding populates pathSummary with hopCount representing path length, estimatedWeight from GraphService cost calculations, and circuitsPreview array providing abbreviated path visualization. Path length exceeding three circuits triggers abbreviated preview generation displaying first circuit, second circuit, ellipsis indication, and final circuit maintaining UI readability while conveying path complexity.
Point Machine Movement Analysis: analyzeRequiredPMMovements method processes discovered paths identifying point machines requiring position changes for route establishment. Direction-dependent target side determination maps UP direction to LEFT side and DOWN direction to RIGHT side establishing point machine positioning requirements for bidirectional track configurations.
Track Circuit Clearance Assessment: checkPathClearance method validates path feasibility through individual circuit occupancy verification using getTrackCircuitOccupancy database queries. Occupied circuits trigger OCCUPIED blocking with specific circuit identification, while resource reservation conflicts detected through isCircuitReserved queries generate RESERVED blocking conditions with affected resource identification.
Point Machine Settability Verification: Point machine operation feasibility assessment through isPointMachineSettable method validates mechanical availability for required position changes. Settable machines populate requiredPMActions array with machine identification, current position status, and target position requirements, while non-settable machines generate LOCKED_PM blocking conditions with specific machine failure identification.
Resource Conflict Detection Framework
Active Route Conflict Analysis: isCircuitReserved method constructs complex database queries joining resource_locks and route_assignments tables identifying circuit reservations by active routes. Query execution with TRACK_CIRCUIT resource type filtering and circuit identification parameters returns boolean results indicating reservation conflicts preventing candidate route establishment.
Point Machine Lock Assessment: Point machine settability determination evaluates machine operating status, lock conditions, transition states, failure modes, and maintenance flags preventing unauthorized operations. Database queries against point_machines table extract current_position_id, operating_status enumeration, is_locked boolean flags, and related status indicators determining operational feasibility.
Resource Lock Type Evaluation: Database queries incorporate lock type differentiation including EXCLUSIVE locks preventing any concurrent access, SHARED locks permitting compatible operations, OVERLAP locks restricting safety margin areas, and EMERGENCY locks overriding standard conflict resolution enabling critical operations during system emergencies.
Reachability Classification Logic
Multi-State Reachability Determination: Candidate classification utilizes hierarchical decision tree evaluating pathfinding success, clearance analysis results, and point machine requirements. REACHABLE_CLEAR designation requires successful pathfinding with complete clearance and no point machine movements, REACHABLE_REQUIRES_PM classification indicates viable paths with necessary point machine operations, BLOCKED categorization encompasses pathfinding failures, occupancy conflicts, resource reservations, and mechanical unavailability.
Failure Reason Categorization: Blocked candidates receive specific failure identification including NO_PATH_FOUND with pathfinding error details, OCCUPIED with circuit identification arrays, RESERVED with conflicting route information, LOCKED_PM with machine status details, and INCOMPLETE_TOPOLOGY with missing connectivity descriptions enabling operator understanding of blocking conditions.
Path Invalidation Handling: Pathfinding failures trigger path metrics invalidation with hopCount set to negative one, estimatedWeight reset to negative one, and circuitsPreview cleared preventing misleading path information display while maintaining candidate object structural integrity for consistent result processing.
Results Processing and Optimization
Candidate Sorting Algorithm: Multi-criteria sorting algorithm prioritizes candidates using hierarchical comparison functions. Primary sorting by reachability priority assigns REACHABLE_CLEAR precedence over REACHABLE_REQUIRES_PM over BLOCKED classifications. Secondary sorting utilizes hop count minimization favoring shorter paths, while tertiary sorting applies estimated weight optimization promoting cost-effective routes.
Statistical Aggregation: Scan completion triggers comprehensive statistical analysis counting reachableClear candidates, reachableRequiresPM alternatives, blocked destinations, and invalid path entries. Statistical data supports operational analysis, system performance monitoring, and user interface display optimization through categorized result presentation.
Result Filtering Mechanisms: includeBlocked parameter controls blocked candidate inclusion through std::remove_if algorithm application with lambda function predicate testing reachability classification. Filtered results maintain statistical accuracy while reducing UI complexity for operators focusing on viable route alternatives.
Comprehensive Result Formatting
Structured Data Organization: formatScanResults method transforms candidate objects into hierarchical QVariantMap structure with reachable_clear array, reachable_requires_pm array, and blocked array organization. Individual candidate formatting includes dest_signal_id, display_name, direction, reachability status, blocked_reason explanations, path_summary objects, required_pm_actions arrays, and conflicts identification arrays.
Performance Metrics Integration: Result objects incorporate scan_time_ms from QElapsedTimer measurements, source_signal_id reference, direction specification, total_candidates count, and summary_stats object containing categorical counts enabling performance analysis, system optimization, and operator workflow enhancement through comprehensive scanning statistics.
Point Machine Action Serialization: RequiredPMAction objects serialize into QVariantMap structures containing machine_id identification, current_position status, and target_position requirements enabling UI display of necessary mechanical operations and operator planning for route establishment procedures requiring infrastructure modifications.
User Interface Integration Framework
QML Component Interaction: Route Assignment Dialog component invokes scanning through globalRouteAssignmentService.scanDestinationSignals with sourceSignalId extraction, AUTO direction specification, and true includeBlocked parameter enabling comprehensive destination analysis. Scanning state management through isScanning boolean flag controls UI responsiveness and progress indication during asynchronous processing.
Result Processing and Display: Successful scanning populates scanResults property triggering UI updates through property binding mechanisms. Result categorization enables tabbed display with reachable_clear, reachable_requires_pm, and blocked sections providing organized candidate presentation. Individual candidate selection through selectDestination method updates selectedDestSignalId property enabling route assignment workflow continuation.
Real-Time Feedback Integration: Console logging provides detailed scanning feedback including candidate counts by category, processing duration, and failure conditions enabling development analysis and operational monitoring. Error handling displays specific failure messages through console.error outputs supporting troubleshooting and system diagnostics during scanning operations.